% Vietnamese translation for OI Public Library
% Source-PDF: 数学 MATHEMATICS/简单数理逻辑及应用 - 李恺威 .pdf
\documentclass[11pt]{article}
\usepackage{oipl-vn}

\newcommand{\True}{\mathrm{T}}
\newcommand{\False}{\mathrm{F}}
\newcommand{\CNF}{\mathrm{CNF}}
\newcommand{\DNF}{\mathrm{DNF}}

\title{Logic toán học đơn giản và ứng dụng}
\author{Li Kaiwei\\Khoa Khoa học và Công nghệ Máy tính, Đại học Thanh Hoa}
\date{}

\begin{document}
\maketitle

\origtitle{Simple mathematical logic and its applications}
\sourcepath{Mathematics / Simple mathematical logic and applications - Li Kaiwei}

\begin{abstract}
Bài giảng giới thiệu các khái niệm cơ bản của logic mệnh đề: mệnh đề, biến
mệnh đề, liên kết logic, công thức hợp lệ, công thức tương đương và dạng
chuẩn. Từ đó tài liệu chuyển sang bài toán SAT, trường hợp đặc biệt 2-SAT,
thuật toán DPLL, rồi mở rộng sang SMT với các miền như số nguyên, số thực,
mảng và bit-vector. Phần cuối minh họa cách giải hệ tuyến tính trên
bit-vector và các ứng dụng trong giải phương trình, sinh kiểm thử và phát
hiện lỗi chương trình.
\end{abstract}

\tableofcontents

\section{Logic mệnh đề}

\subsection{Mệnh đề và biến mệnh đề}

Một \term{mệnh đề} là một câu trần thuật có giá trị chân lý xác định: hoặc
đúng, hoặc sai. Ví dụ ``tuyết màu trắng'' là một mệnh đề đúng; câu tự tham
chiếu kiểu ``điều tôi đang nói là sai'' cho thấy nếu phát biểu không có giá
trị chân lý ổn định thì cần xử lý cẩn thận hơn.

Khi hình thức hóa, ta dùng ký hiệu như \(P,Q,R\) để biểu diễn các mệnh đề.
Một \term{biến mệnh đề} có thể nhận giá trị \(\True\) hoặc \(\False\). Một
mệnh đề phức hợp có thể phân tách thành nhiều mệnh đề đơn. Ví dụ
\[
  P:\quad \text{``tuyết màu trắng và }1+1=2\text{''}
\]
có thể tách thành
\[
  R:\text{``tuyết màu trắng''},\qquad
  S:\text{``}1+1=2\text{''}.
\]

\subsection{Liên kết logic}

Các liên kết thường dùng gồm phủ định \(\neg\), hội \(\land\), tuyển
\(\lor\), kéo theo \(\to\), và tương đương \(\leftrightarrow\). Bảng chân lý
cơ bản:
\[
\begin{array}{c|c}
p & \neg p\\ \hline
\False & \True\\
\True & \False
\end{array}
\qquad
\begin{array}{c|c|c|c|c|c}
p&q&p\land q&p\lor q&p\to q&p\leftrightarrow q\\ \hline
\False&\False&\False&\False&\True&\True\\
\False&\True&\False&\True&\True&\False\\
\True&\False&\False&\True&\False&\False\\
\True&\True&\True&\True&\True&\True
\end{array}
\]

Trong công thức \(p\to q\), giá trị sai chỉ xuất hiện khi tiền đề đúng nhưng
kết luận sai. Điều này phản ánh cách logic hình thức hóa ``nếu \(p\) thì
\(q\)'', không nhất thiết là quan hệ nhân quả trong đời sống.

\subsection{Công thức hợp lệ}

Một \term{công thức hợp lệ} là chuỗi ký hiệu được tạo từ biến mệnh đề và các
liên kết theo quy tắc cú pháp:
\begin{enumerate}
  \item Mỗi biến mệnh đề là một công thức hợp lệ.
  \item Nếu \(A\) là công thức hợp lệ thì \(\neg A\) cũng là công thức hợp lệ.
  \item Nếu \(A,B\) là công thức hợp lệ thì
  \((A\land B)\), \((A\lor B)\), \((A\to B)\), và
  \((A\leftrightarrow B)\) đều là công thức hợp lệ.
  \item Chỉ các chuỗi thu được sau hữu hạn lần dùng ba quy tắc trên mới là
  công thức hợp lệ.
\end{enumerate}

Ví dụ:
\[
  (p\land(p\to q))\to q
\]
là một công thức hợp lệ. Một biểu thức lập trình như
\texttt{if A then B else C} cũng có thể được biểu diễn bằng công thức logic,
chẳng hạn
\[
  (A\land B)\lor(\neg A\land C).
\]

\subsection{Phân loại công thức}

Với một diễn giải \(I\), mỗi biến mệnh đề được gán giá trị đúng hoặc sai.
Một công thức được gọi là:
\begin{itemize}
  \item \term{hằng đúng} nếu nó đúng dưới mọi diễn giải;
  \item \term{khả thỏa} nếu nó đúng dưới ít nhất một diễn giải;
  \item \term{mâu thuẫn} nếu nó sai dưới mọi diễn giải.
\end{itemize}

Các quan hệ cơ bản:
\[
\begin{aligned}
A\text{ hằng đúng} &\Longleftrightarrow \neg A\text{ hằng sai},\\
A\text{ khả thỏa} &\Longleftrightarrow \neg A\text{ không hằng đúng},\\
A\text{ không khả thỏa} &\Longleftrightarrow A\text{ hằng sai}.
\end{aligned}
\]

\section{Tương đương và dạng chuẩn}

\subsection{Công thức tương đương}

Hai công thức \(A,B\) được gọi là \term{tương đương}, ký hiệu \(A=B\), nếu
chúng có cùng giá trị chân lý dưới mọi diễn giải của các biến xuất hiện trong
chúng. Khi đó \(A=B\) khi và chỉ khi \(A\leftrightarrow B\) là hằng đúng.
Ký hiệu \(=\) ở đây không phải là một liên kết logic mới, mà là quan hệ giữa
hai công thức. Quan hệ này có tính phản xạ, đối xứng và bắc cầu.

Một số luật tương đương thường dùng:
\[
\begin{array}{ll}
\text{Phủ định kép:} & \neg\neg P=P,\\[0.2em]
\text{Giao hoán:} & P\lor Q=Q\lor P,\quad P\land Q=Q\land P,\\[0.2em]
\text{Kết hợp:} & (P\lor Q)\lor R=P\lor(Q\lor R),\\
& (P\land Q)\land R=P\land(Q\land R),\\[0.2em]
\text{Phân phối:} & P\lor(Q\land R)=(P\lor Q)\land(P\lor R),\\
& P\land(Q\lor R)=(P\land Q)\lor(P\land R),\\[0.2em]
\text{Lũy đẳng:} & P\lor P=P,\quad P\land P=P,\\[0.2em]
\text{Hấp thụ:} & P\lor(P\land Q)=P,\quad P\land(P\lor Q)=P,\\[0.2em]
\text{De Morgan:} & \neg(P\lor Q)=\neg P\land\neg Q,\\
& \neg(P\land Q)=\neg P\lor\neg Q.
\end{array}
\]
Ngoài ra,
\[
  A\to B=\neg A\lor B,
\]
và
\[
  A\leftrightarrow B=(\neg A\lor B)\land(A\lor\neg B)
  =(A\land B)\lor(\neg A\land\neg B).
\]

\subsection{Từ bảng chân lý đến công thức}

Viết bảng chân lý từ một công thức là dễ. Chiều ngược lại cũng làm được bằng
cách viết từ các dòng đúng hoặc từ các dòng sai. Ví dụ với hai biến \(P,Q\),
giả sử công thức \(A\) đúng ở các dòng \((\False,\False)\),
\((\False,\True)\), \((\True,\True)\), và sai ở
\((\True,\False)\). Từ các dòng đúng, ta viết:
\[
  A=(\neg P\land\neg Q)\lor(\neg P\land Q)\lor(P\land Q).
\]
Từ dòng sai duy nhất, ta có thể viết gọn:
\[
  A=\neg P\lor Q.
\]

\subsection{Dạng chuẩn}

Một \term{literal} là biến mệnh đề \(P\) hoặc phủ định của nó \(\neg P\).
\term{Mệnh đề hội} là hội của một số literal; \term{mệnh đề tuyển} là tuyển
của một số literal.

Một công thức ở dạng \term{tuyển chuẩn} nếu nó là tuyển của các mệnh đề hội:
\[
  A_1\lor A_2\lor\cdots\lor A_n,
\]
trong đó mỗi \(A_i\) là hội của các literal. Một công thức ở dạng
\term{hội chuẩn} nếu nó là hội của các mệnh đề tuyển:
\[
  A_1\land A_2\land\cdots\land A_n,
\]
trong đó mỗi \(A_i\) là tuyển của các literal.

Định lý dạng chuẩn nói rằng mọi công thức mệnh đề đều có một công thức tương
đương ở dạng hội chuẩn và một công thức tương đương ở dạng tuyển chuẩn. Các
luật như \(A\to B=\neg A\lor B\), De Morgan và phân phối là công cụ cơ bản để
biến đổi.

\section{Bài toán SAT}

\subsection{Định nghĩa}

\term{SAT} là bài toán quyết định xem một công thức logic có khả thỏa hay
không. Thông thường công thức được đưa về dạng hội chuẩn:
\[
  A_1\land A_2\land\cdots\land A_n,
\]
trong đó mỗi \(A_i\) là một mệnh đề tuyển
\[
  A_i=(P_{i1}\lor P_{i2}\lor\cdots\lor P_{im}).
\]

\subsection{2-SAT}

2-SAT là trường hợp đặc biệt trong đó mỗi mệnh đề tuyển có không quá hai
literal. Ví dụ:
\[
  (X_0\lor X_2)\land(\neg X_0\lor X_3)\land(X_1\lor\neg X_3).
\]
2-SAT có lời giải đa thức bằng đồ thị kéo theo. Với mỗi biến \(A\), tạo hai
đỉnh \(A\) và \(\neg A\). Với mỗi mệnh đề \((A\lor B)\), dùng tương đương
\[
  A\lor B \equiv (\neg A\to B)\land(\neg B\to A),
\]
nên thêm hai cung
\[
  \neg A\to B,\qquad \neg B\to A.
\]

Nếu trong đồ thị có cả đường đi \(A\to^{*}\neg A\) và
\(\neg A\to^{*}A\), tức \(A\) và \(\neg A\) nằm trong cùng một thành phần
liên thông mạnh, thì công thức vô nghiệm. Ngược lại, thu gọn SCC thành DAG và
chọn giá trị cho từng cặp đỉnh đối ngẫu theo thứ tự topo ngược để có một gán
thỏa.

\subsection{3-SAT và DPLL}

Khi có mệnh đề tuyển chứa ba biến, cách làm trên không còn áp dụng trực tiếp.
3-SAT là bài toán NP-đầy đủ đầu tiên được biết đến, thông qua định lý Cook
năm 1971. Với SAT tổng quát, ta thường phải tìm kiếm có cắt tỉa.

Thuật toán Davis--Putnam--Logemann--Loveland, gọi tắt là \term{DPLL}, được
đưa ra năm 1962 như một cải tiến của thuật toán Davis--Putnam năm 1960. Trong
nhiều thập niên, nó vẫn là lõi của nhiều bộ giải SAT hiệu quả.

Dạng khung:
\begin{verbatim}
function DPLL(Phi):
    if Phi is empty:
        return true
    if Phi contains an empty clause:
        return false

    while Phi has a unit clause l:
        set l to true and simplify Phi

    while Phi has a pure literal x:
        set x to true and simplify Phi

    choose an unassigned variable y
    return DPLL(Phi with y = true)
        or DPLL(Phi with y = false)
\end{verbatim}

Ở đây \(\Phi\) là tập các mệnh đề tuyển, được hiểu là hội của chúng. Hai cắt
tỉa quan trọng là \term{unit propagation}, tức mệnh đề chỉ còn một literal
thì literal đó buộc phải đúng, và \term{pure literal elimination}, tức nếu một
biến chỉ xuất hiện với cùng một dấu thì có thể gán nó để làm đúng tất cả các
mệnh đề chứa nó.

\section{SMT}

\subsection{Từ SAT đến SMT}

SAT chỉ xử lý ràng buộc trên biến Boolean. Trong nhiều bài toán, ràng buộc
không chỉ là tuyển các biến Boolean, mà còn là phương trình hoặc bất phương
trình trên số nguyên, số thực, mảng, hoặc bit-vector. \term{SMT}, viết tắt
của \textit{Satisfiability Modulo Theories}, là bài toán xét khả thỏa của
công thức logic trên một lý thuyết nền cụ thể.

Một số miền thường gặp:
\begin{itemize}
  \item Boolean, tức SAT cổ điển;
  \item số nguyên;
  \item số thực;
  \item mảng;
  \item bit-vector.
\end{itemize}

SMT có hai hướng nghiên cứu lớn. Một hướng gần với toán học: giải ràng buộc
trên số nguyên, số thực, tuyến tính hoặc phi tuyến. Hướng còn lại gần với
máy tính: bit-vector, mảng và các phép toán mô phỏng chính xác hành vi của
chương trình.

\subsection{Ví dụ giải hệ tuyến tính trên bit-vector}

Xét hệ trên bit-vector độ rộng \(3\):
\[
\begin{aligned}
3x+4y+2z&=0,\\
2x+2y+2&=0,\\
4y+2x+2z&=0.
\end{aligned}
\]
Các phép toán được hiểu theo modulo \(2^3=8\). Từ phương trình đầu, vì
\(3^{-1}\equiv3\pmod 8\), có thể biểu diễn
\[
  x=4y+2z.
\]
Thay vào hai phương trình còn lại:
\[
\begin{aligned}
2y+4z+2&=0,\\
4y+6z&=0.
\end{aligned}
\]
Tất cả hệ số đều chẵn nên chia cho \(2\), đồng thời giảm độ rộng còn \(2\)
bit và bỏ bit cao:
\[
\begin{aligned}
y[1:0]+2z[1:0]+1&=0,\\
2y[1:0]+3z[1:0]&=0.
\end{aligned}
\]
Giải \(y[1:0]\) theo \(z[1:0]\):
\[
  y[1:0]=2z[1:0]+3.
\]
Thay vào phương trình còn lại:
\[
  3z[1:0]+2=0,
\]
từ đó thu được
\[
  z[1:0]=2,\qquad y[1:0]=3.
\]
Quay lại độ rộng \(3\), viết \(y=y'\mathbin{@}2\), \(z=z'\mathbin{@}3\) theo
cách ký hiệu nối bit trong slide, và
\[
  x=4y+2z.
\]
Ví dụ này minh họa đặc trưng của bit-vector: phép chia, bỏ bit cao và đồng dư
theo độ rộng đều là một phần của ngữ nghĩa.

\section{Ứng dụng}

\subsection{Giải phương trình}

SMT có thể được dùng trực tiếp để kiểm tra một hệ phương trình hoặc bất
phương trình có nghiệm hay không, đồng thời trả về một nghiệm nếu tồn tại.
Tùy miền, bộ giải sẽ dùng các thủ tục khác nhau: Gaussian elimination cho hệ
tuyến tính, simplex cho ràng buộc tuyến tính trên số thực, hoặc các thủ tục
bit-vector cho phép toán mức máy.

\subsection{Phát hiện lỗi chương trình}

Một ứng dụng quan trọng là quét lỗi chương trình bằng cách biến đường đi thực
thi thành ràng buộc. Ví dụ:
\begin{verbatim}
int two_hop(int x) {
    int a[4] = {3, 0, 2, 1};
    if (x < 0 || x > 3) return -1;
    return a[a[x] - 1];
}
\end{verbatim}
Điều kiện đầu vào bảo đảm \(x\in[0,3]\), nhưng khi \(x=1\) thì
\(a[x]=0\), nên truy cập \(a[a[x]-1]=a[-1]\), vượt ngoài mảng. Một bộ giải SMT
có thể sinh tự động đầu vào \(x=1\) bằng cách mô hình hóa mảng, phép so sánh
và phép truy cập theo đường đi.

\section*{Tài liệu tham khảo}

\begin{itemize}
  \item Shi Chunyi, \textit{Mathematical Logic and Set Theory}, Tsinghua
  University Press, 2000.
  \item Zhao Shuang, \textit{Analysis of 2-SAT solving methods}.
  \item Wikipedia: Boolean satisfiability problem.
  \item Wikipedia: DPLL algorithm.
  \item Wikipedia: Satisfiability Modulo Theories.
  \item Cristian Cadar, Vijay Ganesh, Peter Pawlowski, David Dill, and Dawson
  Engler. \textit{EXE: Automatically generating inputs of death}. ACM CCS,
  2006.
  \item Vijay Ganesh. \textit{Decision Procedures for Bit-vectors, Arrays and
  Integers}. September 2007.
\end{itemize}

\end{document}
